\section{Ordinary arithmetic and the two interpretation composites}
\label{sec:main-ordinary}

The rule-equivalence theorem is only one part of the requirement. This
section states precisely which ordinary first-order theory the calculus
presents and why both interpretation composites preserve the whole
numerical domain.

\subsection{An explicit ordinary realization}

Let $T_I$ use intuitionistic first-order logic with actual equality and
genuine falsity. Its numerical language contains $0$, $S$, and every fixed
finite scalar-iterator description. Formally, at each stage, for an
earlier term $h(u,\vec p)$ add a function symbol $i_h(n,a,\vec p)$ with
\[
 i_h(0,a,\vec p)=a,\qquad
 i_h(Sn,a,\vec p)=h(i_h(n,a,\vec p),\vec p).
\]
The remaining nonlogical axioms or equivalent contextual rules are
\[
 \forall x\,(x=0\lor\exists y\,x=Sy),\qquad
 0=S0\longrightarrow\bot,
\]
and induction for each literal equation $E(n,\vec p)$, with arbitrary
retained assumptions and the usual fresh induction parameter.
Only fixed finite descriptions are added. There is no numerical
quantification over programs and no total universal evaluator.

The ordinary readings are
\[
 (s=t)^\circ=(s=t),\qquad
 (D_{x,y}B)^\circ=\exists x\bigl(B^\circ(x,0)\to
                                      \forall y B^\circ(x,y)\bigr),
\]
\[
 [A;y.B]^\circ=A^\circ\to\forall y B^\circ.
\]
Every primitive rule of $\CE$ is sound under this reading. Conversely,
the equality rules, the ordinary logical rules, coverage, explosion,
and equation induction are derived in $\CE$ by its equivalence with
the reference clause calculus. The explicit falsity bridge matters:
equational computation alone permits a singleton algebra and must not
silently be identified with an inconsistent ordinary theory.

\subsection{What the inherited scalar bootstrap proves}

The appendix supplies the existing internal argument in its necessary
dependency order:
\begin{enumerate}
\item Construct a scalar predecessor $P$ and prove $P0=0$ and
      $P(Sn)=n$ using equation induction and numerical cases.
\item Construct truncated subtraction and a faithful characteristic
      of actual numerical equality. The reflection theorem holds at
      arbitrary elements, including nonstandard elements of a model.
\item Translate each fixed quantifier-free formula to a characteristic
      equation and obtain its full intuitionistic induction instance.
\item Construct pairing and projections by bounded search, then compile
      each fixed general primitive recursion to scalar iteration.
\end{enumerate}
The order avoids assuming pairing or a faithful equality detector while
constructing predecessor. These constructions are inherited rather than
new size-reduction claims of this report.

If $\pi,L,H$ are the resulting pairing and projections, a target recursion
\[
 f(0,\vec p)=g(\vec p),\qquad
 f(Sn,\vec p)=h(n,f(n,\vec p),\vec p)
\]
is compiled as
\[
 \begin{split}
 V(n,\vec p)&=I\bigl(n,\pi(0,g^*(\vec p));
        u.\pi(SL(u),h^*(L(u),H(u),\vec p))\bigr),\\
 f^*(n,\vec p)&=H(V(n,\vec p)).
 \end{split}
\]
Internal equation induction proves $L(V(n,\vec p))=n$, so the projection
equations prove both defining equations for $f^*$. This is a finite
external compilation on descriptions.

\subsection{Bi-interpretability, not just agreement on numerals}

Interpret a scalar iterator in ordinary $\ipra$ by primitive recursion
whose update ignores its counter. Interpret every $\ipra$ description
back in $T_I$ by the compiler above. In both directions the domain
formula is $x=x$ and numerical equality remains actual equality.

For two fixed iterators with equal seeds and provably equal updates,
equation induction proves equality of all their values. The analogous
argument proves uniqueness of each ordinary primitive recursion.
Induction externally on a finite description now shows that its double
translation satisfies its original defining equations and hence is
internally equal to the original function. This also checks coherence
when substitution changes an update's syntactic description.

Consequently the graph $x=y$ is internally a total bijection that
preserves every numerical operation in each composite interpretation.
For formulas, the crucial native composite is
\[
 D_{x,y}B\ \dashv\vdash\
 \exists_*x\bigl(B(x,0)\Rightarrow_*\forall_*y B(x,y)\bigr).
\]
The forward direction opens its clause, assumes the base, and derives
every instance; the backward direction uses the supplied universal in
case introduction. The remaining formula composites follow by external
structural induction through the explicitly derived logical macros.
Thus $T_I$, and therefore the ordinary realization of $\CE$, is
bi-interpretable and equiconsistent with full first-order $\ipra$.

This is a written mathematical proof using the supplied internal
arithmetic development. The executable tests below concern the new
checker implementation and finite proof translations; they are not a
replacement for the interpretation proof and are not advertised as a
new Lean or Rocq formalization.
